\section{Appendix}
\subsection{Proposition Proof}
\label{appx:proof}

\begin{proof}
\begin{equation} 
\begin{split}
    &\mathbb{E}(z_i|\mathbf{x}_i,\mathit{G})=\mathbb{E}(y_i^{obs}\cdot \frac{t_i-p}{p(1-p)}|\mathbf{x}_i,\mathit{G})  \\ 
    &=\mathbb{E}(t_iy_i^{obs}\cdot \frac{t_i-p}{p(1-p)}+(1-t_i)y_i^{obs}\cdot \frac{t_i-p}{p(1-p)}|\mathbf{x}_i,\mathit{G})
\end{split}
\label{eqn:proof1}
\end{equation}
Given $t_i\in\{0,1\}$, we have $t_i^2=t_i$ and $(1-t_i)^2=(1-t_i)$. Simultaneously with the definition in Eq. \eqref{eqn:y_obs}, we have:
\begin{equation}
\left\{
             \begin{array}{lr}
             t_iy_i^{obs}=t_i^2y_i(1) + t_i(1 - t_i)y_i(0)=t_i^2 y_i(1)=t_iy_i(1) \nonumber \\ 
(1-t_i)y_i^{obs}=t_i(1-t_i)y_i(1) + (1 - t_i)^2y_i(0)=(1-t_i)y_i(0) . \nonumber
             \end{array}
\right.
\end{equation}

Then we can rewire Eq. \ref{eqn:proof1} as \footnote{Due to the limit of space, $t_i^j$ represents $t_i=j$ here.}: 
\begin{equation} 
\begin{split}
    &\mathbb{E}(t_iy_i^{obs}\cdot \frac{t_i-p}{p(1-p)}+(1-t_i)y_i^{obs}\cdot \frac{t_i-p}{p(1-p)}|\mathbf{x}_i,\mathit{G}) \\
    &=\mathbb{E}(t_iy_i(1)\frac{t_i-p}{p(1-p)}+(1-t_i)y_i(0) \frac{t_i-p}{p(1-p)}|\mathbf{x}_i,\mathit{G},t_i^1)*p(t_i^1|\mathbf{x}_i,\mathit{G})\\
    &+\mathbb{E}(t_iy_i(1)\frac{t_i-p}{p(1-p)}+(1-t_i)y_i(0) \frac{t_i-p}{p(1-p)}|\mathbf{x}_i,\mathit{G},t_i^0)*p(t_i^0|\mathbf{x}_i,\mathit{G}) \\
    &=\mathbb{E}(\frac{y_i(1)}{p}|\mathbf{x}_i,\mathit{G})*p(t_i=1|\mathbf{x}_i,\mathit{G})-\mathbb{E}(\frac{y_i(0)}{1-p}|\mathbf{x}_i,\mathit{G})*p(t_i=0|\mathbf{x}_i,\mathit{G})\\
    &=\mathbb{E}(y_i(1)|\mathbf{x}_i,\mathit{G})\frac{p(t_i=1|\mathbf{x}_i,\mathit{G})}{p}-\mathbb{E}(y_i(0)|\mathbf{x}_i,\mathit{G})\frac{p(t_i=0|\mathbf{x}_i,\mathit{G})}{1-p}\\
    &=\mathbb{E}(y_i(1)|\mathbf{x}_i,\mathit{G})-\mathbb{E}(y_i(0)|\mathbf{x}_i,\mathit{G}) \\
    &=\hat{\tau}_i \nonumber
\end{split}
\end{equation}
\end{proof}

\subsection{Pseudo Code and Complexity Analysis}
\label{appx:code}
The pseudo-code of GNUM-CT is given in Algorithm \ref{alg2} and GNUM-PL in Algorithm \ref{alg1}. Our proposed method can be trained by an end-to-end back-propagation, and thus we can use gradient descent to optimize the model.

For each sample, it will go through the GNUM layer first and then perform the uplift prediction. The complexity of the whole process is
$O\left(M\sum_{i=0}^L f_l+N \sum_{i=1}^L f_{i-1}f_i \right)$ where $N$ denotes the total number of users, including the users in the treatment and control group, $M = \left| \mathit{E}\right|$ is the number of relationships, $L$ is number of hidden layers and $f_l$ is the dimensionality of the $l^{th}$ hidden layer. In practice, $L$ and $f_l$ are often bounded within a small constant. The proposed method is linear to the number of users and number of relationships in the dataset respectively. Therefore, the overall model is scalable. 

$O(Nd^2_{max}S_{max})$, where $N$ denotes the total number of users, including the users in the treatment and control group, $d_{max}$ denotes the maximum number of dimensionality among different layers and $S_{max}$ denotes in the social graph the maximum degree among all the users. In practice, $S_{max}$ is often bounded within a constant. For example, in social networks like Facebook, the number of friends a user can have has an upper bound. 

\begin{algorithm}[t]\caption{Graph Neural Network for Uplift Modeling (GNUM) using Class-Transformed Target} \label{alg2}
	\begin{algorithmic}[1]
		\REQUIRE Set of user features, treatment assignment and observed outcome $\{(\mathbf{x}_i,t_i,y_i)\}_{i=1}^N$, social graph $\mathit{G}=(\mathit{V},\mathit{E})$ and number of layers $L$, 
		\ENSURE Prediction of user uplift and GNUM parameters $\mathbf{\Theta}$.
		\STATE Initialize all parameters $\mathbf{\Theta}$ and using $\mathbf{X}$ as $\mathbf{H}^{(0)}$
		\WHILE {$\mathcal{L}$ does not converge}
		\FOR{i $\leftarrow$ 1 to L}
		\STATE Calculate graph-based representation $\mathbf{H}^{(l)}$.
		\ENDFOR
		\STATE Calculate  $\hat{y}_i^{pl1}$ and $\hat{y}_i^{pl2}$  from learned representations $\mathbf{H}^{(L)}$ using Eq. \eqref{eqn:partial_resp}.
		\STATE Calculate the class-transformed target $z_i$ using Eq. \eqref{eq_zi}.
		\STATE Calculate the loss $\mathcal{L}_{CT}$ using Eq. \eqref{eq:loss_ct}.
		\STATE Update $\mathbf{\Theta}$ using back-propagation.
		\ENDWHILE
		\STATE Using the trained model parameters to estimate the user uplift as: $\hat{\tau}_i=\mathbb{E}(z_i|\mathbf{x}_i, \mathit{G})$.
	\end{algorithmic}
\end{algorithm}

\begin{algorithm}[t]\caption{Graph Neural Network for Uplift Modeling (GNUM) using Partial Label Learning} \label{alg1}
	\begin{algorithmic}[1]
		\REQUIRE Set of user features, treatment assignment and observed outcome $\{(\mathbf{x}_i,t_i,y_i)\}_{i=1}^N$, social graph $\mathit{G}=(\mathit{V},\mathit{E})$ and number of layers $L$, 
		\ENSURE Prediction of user uplift and GNUM parameters $\mathbf{\Theta}$.
		\STATE Initialize all parameters $\mathbf{\Theta}$ and using $\mathbf{X}$ as $\mathbf{H}^{(0)}$
		\WHILE {$\mathcal{L}$ does not converge}
		\FOR{i $\leftarrow$ 1 to L}
		\STATE Calculate graph-based representation $\mathbf{H}^{(l)}$.
		\ENDFOR
		\STATE Calculate  $\hat{y}_i^{pl1}$ and $\hat{y}_i^{pl2}$  from learned representations $\mathbf{H}^{(L)}$ using Eq. \eqref{eqn:partial_resp}.
		\STATE Construct the sample's partial labels based on $y_i$ and $t_i$.
		\STATE Calculate the loss $\mathcal{L}$ using Eq. \eqref{eq:gnum_loss}.
		\STATE Update $\mathbf{\Theta}$ using back-propagation.
		\ENDWHILE
		\STATE Using the trained model parameters to estimate the user uplift as: $\hat{\tau}_i = 1 - P(S_i=[1, 0, 0]|\mathbf{x}_i,\mathit{G}) - P(S_i=[0, 0, 1]|\mathbf{x}_i,\mathit{G})$
	\end{algorithmic}
\end{algorithm}



\subsection{More Experimental Details}
\subsubsection{More Details about Dataset}
\label{appn:dataset}
Blogcatalog is a social blog directory that manages bloggers and their blogs. The dateset contains 
the features of bloggers and the social relationships between bloggers listed on the BlogCatalog website. The features are bag-of-words representations of keywords in bloggers’ descriptions. 
We follow the assumptions and procedures
of synthesizing the outcomes and treatments assignments in \cite{guo2020learning}, in which the outcomes represent the opinions of readers on each
blogger and the treatments represent whether contents created by a blogger receive more views on mobile devices or desktops.
In the treatment group, the blogger’s blogs are read more on mobile devices. In the control group, the blogger’s blogs are read more on desktops. 
We assume that the social relationships of bloggers can causally affect their treatment assignments and readers’ opinions of them. 
We trained the LDA topic model to get the device
preference of the readers of the i-th blogger’s content as:

\begin{center}
$
\begin{aligned}
\operatorname{Pr}(t&\left.=1 \mid \mathbf{x}_i, \mathrm{~A}\right)=\frac{\exp \left(p_1^i\right)}{\exp \left(p_1^i\right)+\exp \left(p_0^i\right)} ; \\
p_1^i &=\kappa_1 r\left(\mathbf{x}_i\right)^T r_1^c+\kappa_2 \sum_{j \in \mathcal{N}(i)} r\left(\mathbf{x}_j\right)^T r_1^c \\
&=\kappa_1 r\left(\mathbf{x}_i\right)^T r_1^c+\kappa_2\left(\mathrm{~A} r\left(\mathbf{x}_j\right)\right)^T r_1^c ; \\
p_0^i &=\kappa_1 r\left(\mathbf{x}_i\right)^T r_0^c+\kappa_2 \sum_{j \in \mathcal{N}(i)} r\left(\mathbf{x}_j\right)^T r_0^c \\
&=\kappa_1 r\left(\mathbf{x}_i\right)^T r_0^c+\kappa_2\left(\mathrm{~A} r\left(\mathbf{x}_j\right)\right)^T r_0^c,
\end{aligned}
$
\end{center}

where $\kappa_1, \kappa_2 \geq 0$ signifies the magnitude of the confounding bias resulting from a blogger’s topics and her neighbors’ topics, respectively.  $\kappa_1 = 0, \kappa_2 = 0$ means
the treatment assignment is random and there is no selection bias,
and greater $\kappa_1, \kappa_2$ means larger selection bias. The factual outcome and the counterfactual outcome of the i-th blogger are simulated as:

\begin{center}
$
\begin{aligned}
y^F\left(\mathbf{x}_i\right) &=y_i=C\left(p_0^i+t_i p_1^i\right)+\epsilon \\
y^{C F}\left(\mathbf{x}_i\right) &=C\left[p_0^i+\left(1-t_i\right) p_1^i\right]+\epsilon,
\end{aligned}
$
\end{center}
where $C$ is a scaling factor and the noise is sampled as $\epsilon \sim \mathcal{N}(0,1)$. In this work, we set $C=5, \kappa_1=10, \kappa_2 \in\{0.5,2\}$.

\subsubsection{Details about the comparing methods}
\label{appn:baseline}
The detail descriptions of the comparing methods can be found here. 
\begin{itemize}
	\item{Two-Model: This method infers the labels in the treatment group and the control group respectively. }
	\item{CTM: The Class-Transformation model (CTM) creates the new target variable to estimate the uplift.}
	\item{Uplift-RF}: This model modifies existing random forest algorithms to directly infer the uplift.
	\item{DML}: Double Machine Learning (DML) method uses the Neyman orthogonal score and cross-fitting to construct the uplift estimator.
	\item{DRL}: Doubly Robust Learning (DRL) method combines the error imputation and the inverse propensity score estimator to address the bias problem.
	\item{NetDeconf}: It uses the graph to minimize confounding bias in the task of estimating treatment effects.
	\item{GNUM-GCN/GAT}: To show the effectiveness of partial label learning, we replace the graph representation learning method with GCN and GAT respectively for comparison.
\end{itemize}


\subsubsection{Parameter Settings}
\label{appn:para}
We adopt the deep neural networks \cite{liu2017survey} with two $64$-unit hidden layers as the building block for the Two-Model and CTM methods. For uplift random forest, we set the number of estimators as $50$ and the depth of the tree as $5$. For NetDeconf and our proposed method, we use two layers of GNN with units of $128-64$. All the weight matrices are initialized using Xavier initialization \cite{glorot2010understanding}.  We train the model for $5$ epochs with a learning rate of $0.0001$ and batch size of $256$. Note that we use grid search to get the best hyper-parameters. For each dataset, we randomly sampled 70\% of the users as the training set, 10\% as the validation set, and 20\% as the test set.  The models are trained on a cluster of 10 Dual-CPU servers with AGL \cite{zhang2020agl} framework. For the largest dataset, containing 573.8K nodes and 2.6M edges, the proposed method took about 17 minutes to train on a cluster of 10 Dual-CPU servers. We run each algorithm $5$ times and report the average result. 

